\documentclass{pma-bthesis}
\begin{document}
\maketitlepage{
    theme = {Тема Вашої дипломної роботи},
    group = {КМ-ХХ},
    author = {Ваші ПІБ},
    supervisorTitles = {регалії керівника роботи},
    supervisor = {ПІБ},
    normConsultantTitles = {регалії консультанта з нормоконтролю},
    normConsultant = {ПІБ},
% за відсутності консультанта приберіть/закоментуйте наступні три рядки
    consultantIntro = {Консультант з \dots}, % розділи, "технічна частина" тощо
    consultantTitles = {регалії консультанта},
    consultant = {ПІБ},
    reviewerTitles = {регалії рецензента},
    reviewer = {ПІБ},
    year = {2077}
}

% завдання на дипломну роботу (впишіть відповідні значення)
\begin{assignment}{Ваші ПІБ у давальному відмінку}
        {Ім'я ПРІЗВИЩЕ (Ваші, для підписів)}
        {Ім'я ПРІЗВИЩЕ (керівника, для підписів)}
    % у першому пункті потрібно вказати тему дипломної роботи, прізвище, ім’я, по батькові (повністю), науковий ступінь, учене звання керівника роботи, а також реквізити наказу, яким їх затверджено.
    \item Тема роботи: \inquot{\theme}, керівник роботи \supervisor{}, \supervisorTitles{}, затверджені наказом по університету від \dateplaceholder{} №~\_\_\_\_-\_.

    % у другому пункті потрібно вказати термін подання переплетеної роботи на кафедру.
    \item Термін подання студентом роботи: \dateplaceholder{}

    % у пункті «Вихідні дані до роботи» потрібно вказати кількісні/якісні показники щодо умов, засобів та методів, які характеризують спрямованість дослідження, конкретизують методику розв'язання задачі та проведення експерименту
    \item Вихідні дані до роботи: \dots

    % у пункті «Зміст роботи» потрібно вказати завдання з окремих частин дипломної роботи. Формулювати ці завдання потрібно в наказовому способі, починаючи зі слів «розробити», «виконати», «проаналізувати», «провести» тощо
    \item Зміст роботи: \dots

    % у пункті «Перелік ілюстративного матеріалу» потрібно вказати назви ключових для дипломної роботи рисунків та слайдів презентації, із якою студент виступатиме під час захисту дипломної роботи
    \item Перелік ілюстративного матеріалу: \dots
    % у пункті «Консультанти розділів роботи» потрібно в окремих рядках таблиці вказати назви розділів, а також прізвища, ініціали та посади консультантів цих розділів (за наявності, у протилежному випадку цей пункт можна пропустити). Варто зазначити, що консультант із нормоконтролю є консультантом усієї роботи, а не окремих розділів, тому в цьому пункті його зазначати непотрібно
    \item
    \begin{consultants}
        Назва розділу & Прізвище та ініціали, посада &&\\\hline
        \dots & \dots &&
    \end{consultants}

    % у пункті «Дата видачі завдання» потрібно вказати дату, коли студенту було видано завдання на дипломну роботу.
    \item Дата видачі завдання: \dateplaceholder{}

    \vspace{1.5em}
    \begin{calenderplan}
        % Впишіть склад календарного плану
        1. & Збір та опрацювання літератури за тематикою дипломної роботи & \dots &\\\hline
        2. & \dots & \dots &
    \end{calenderplan}
\end{assignment}

% «Анотація» повинна стисло відображати загальну характеристику та основний зміст дипломної роботи. Обсяг «Анотації» повинен становити не менше за 650 знаків (із пробілами).
\annotation{Анотація}
% «Анотацію» потрібно виконувати згідно з такою структурою:
% - відомості про обсяг дипломної роботи (БЕЗ додатків), кількість рисунків, таблиць, додатків і бібліографічних найменувань за «Переліком посилань»
Дипломну роботу виконано на \totpages{аркуші}{аркушах}{аркушах}, вона містить \totapp{додаток}{додатки}{додатків} та перелік посилань на використані джерела з \totcit{найменування}{найменувань}{найменувань}. У роботі наведено \totfig{рисунок}{рисунки}{рисунків} та \tottab{таблицю}{таблиці}{таблиць}.

% - мета роботи
% - методи, розглянуті в роботі, критерії, за якими здійснювався вибір методу(-ів) для роботи, метод(-и), вибраний(-і) для використання в роботі
% - основні результати, яких було досягнуто в процесі виконання роботи. При цьому потрібно використовувати безособові форми дієслова доконаного способу минулого часу, наприклад, «виконано», «реалізовано», «проведено», «спроектовано» тощо
% - відомості щодо апробації результатів дипломної роботи (виступи на наукових конференціях, публікації в наукових журналах, довідки про впровадження тощо) чи рекомендації щодо використання результатів
% - ключові слова (від 5 до 20 позицій)

\annotation{Annotation}
% «Abstract» (анотація англійською мовою) повинна бути перекладом анотації українською мовою. Обсяг «Abstract» повинен становити не менше за 650 знаків (із пробілами).
This thesis is submitted on \totpages{page}{pages}{pages}, it has \totapp{appendix}{appendices}{appendices} and \totcit{reference}{references}{references}. There are \totfig{figure}{figures}{figures} and \tottab{table}{tables}{tables} in this document.

\tableofcontents

% «Вступ» має відображати актуальність роботи на основі стислого огляду сучасного стану проблемної області.
\intro

% Першим розділом роботи *обов'язково* повинен бути розділ «Постановка задачі», який повинен містити мету роботи, завдання, які потрібно розв'язати для її досягнення, вимоги до розроблюваних у роботі математичних методів та програмних засобів.
\chapter{Постановка задачі}

% У наступних розділах дипломної роботи потрібно подати:
% - огляд актуальних теоретичних та практичних рішень за тематикою роботи, критерії вибору методу(-ів) для використання в роботі, порівняльний аналіз розглянутих рішень за відповідними критеріями
% - опис побудованої математичної моделі та методу(-ів) розв'язання поставленої задачі
% - опис програмних засобів розв'язання поставленої задачі, у т. ч. розроблених алгоритмів
% - опис випробування розроблених програмних засобів
\chapter{Другий розділ}
Посилання додаються таким чином~\cite{Aho1992,Laming1968,NML}.

Математична формула у тексті: $x+\frac{x}{y} - y^{10} + y^20$, порівняйте степінь в $\{\dots\}$ та без. Формула окремим блоком з нумерацією:
\begin{equation}
    a = \frac{\partial b}{\partial{} y}\times\left.\sum_{i=0}^{\infty}\frac{1}{x^i}\right| _{x\rightarrow7} + 280\sin^3(\gamma)\left(\frac{\tanh(130^\circ - 42\pi n)}{\cosh(\alpha - \beta^2 + \sigma)}+\int_0^5 \frac{x^3 dx}{x + x^5} \right),
\end{equation}
\begin{eqdesc}
    \item[$a$] --- щось;
    \item[$b$] --- вже щось інше;
    \item[\dots] --- \dots.
\end{eqdesc}

Тестова таблиця (див. табл~\ref{tab:testing}).
\begin{table}{|l|r|}{Назва таблиці.}{tab:testing}{\hline a & b\\\hline}
    c & d
\end{table}

% Кожний розділ роботи, окрім «Постановки задачі», повинен закінчуватися підрозділом «Висновки до розділу».
\section{Висновки до розділу}

\chapter{Ще один розділ}

% У «Висновках» потрібно викласти найважливіші результати, одержані в дипломній роботі, потрібно зазначити якісні та кількісні показники відповідних результатів.
\conclusions

% «Перелік посилань» потрібно оформлювати за вимогами, викладеними на сторінці «Вимоги до оформлення документації».
\bibliographystyle{udstu2006}
\bibliography{bibliography}

% У дипломній роботі, серед інших, повинні бути такі обов'язкові додатки:
\appendix

% першим додатком повинен бути додаток «Лістинги програм», до якого потрібно винести лістинги програмного коду реалізації описаних у роботі математичних моделей, методів та алгоритмів. Код графічного інтерфейсу користувача можна не наводити. Не менше 30% рядків лістингів повинні складати коментарі. Обсяг лістингів повинен складати 15–40 сторінок. Варто пам'ятати, що обсяг лістингів не впливає на обсяг дипломної роботи!
\chapter{Лістинги програм}
\lstinputlisting[language=Python, caption=main.py]{main.py} %< таким чином можна включити текст файлу до лістингів; TODO \srcset

% Додаткову інформацію: доведення теорем, матеріали до роботи --- потрібно також подавати у додатках
\chapter{Додаткові матеріали}

% передостаннім додатком повинен бути додаток «Ілюстративний матеріал», до якого потрібно винести слайди презентації, із якою студент виступатиме під час захисту дипломної роботи
\chapter{Ілюстративний матеріал}

% останнім додатком, за наявності відповідних матеріалів, повинен бути додаток «Довідки про впровадження результатів роботи», до якого потрібно винести скан-копії довідок (актів) про впровадження результатів дипломної роботи. У випадку відсутності таких матеріалів цей додаток можна не наводити
\chapter{Довідки про впровадження результатів роботи}

\end{document}